\section*{第 \ref{ch:g11:polarity-and-cohesion} 章——电负性、极性与分子间作用力}

\begin{solution}{exo:g11:polarity-and-cohesion:1}
\ce{H-Cl}：氯（3.16 对 2.20）。\ce{O-H}：氧（3.44）。
\ce{C-O}：氧（3.44 对 2.55）。
\end{solution}

\begin{solution}{exo:g11:polarity-and-cohesion:2}
\ce{H-H}：0，非极性。\ce{C-H}：0.35，视为非极性。\ce{O-H}：
1.24，有极性。\ce{N-H}：0.84，有极性。\ce{Cl-Cl}：0，非极性。\ce{C-Cl}：
0.61，有极性。
\end{solution}

\begin{solution}{exo:g11:polarity-and-cohesion:3}
氮（3.04 > 2.19）与氧（3.44 > 2.58）的\omterm{def:g11:polarity-and-cohesion:electronegativity}{电负性}
自上而下减小。氮（3.04 > 2.55）与镁
（1.31 > 0.93）：沿周期从左到右增大。
\end{solution}

\begin{solution}{exo:g11:polarity-and-cohesion:4}
氨和水：二者都有与氮或氧相连的氢\omterm{def:g7:atoms-and-molecules:atom}{原子}，
且该\omterm{def:g7:atoms-and-molecules:atom}{原子}上有\omterm{def:g10:lewis-and-shape:lone-pair}{孤电子对}。甲烷中没有与 N、O 或 F 相连的氢；
氯化氢中的氢与氯相连，氯
不是定义中三种\omterm{def:g7:atoms-and-molecules:atom}{原子}之一。
\end{solution}

\begin{solution}{exo:g11:polarity-and-cohesion:5}
是的：所有\omterm{def:g7:atoms-and-molecules:molecule}{分子}之间都存在范德华相互作用。但甲烷
\omterm{def:g7:atoms-and-molecules:molecule}{分子}较小（10 个\omterm{def:g9:inside-the-atom:nucleus}{电子}），因此这些相互作用很弱，
在很低的温度下即被破坏：甲烷的沸点约为
\qty{-162}{\celsius}。
\end{solution}

\begin{solution}{exo:g11:polarity-and-cohesion:6}
\ce{CH2Cl2}：极性（两个极性 \ce{C-Cl} 键与两个几乎非极性的
\ce{C-H} 键，无对称性抵消它们）。\ce{CCl4}：非极性（四个
相同的键位于正四面体的顶点）。\ce{NH3}：极性
（三角锥形）。\ce{CO2}：非极性（直线形，两键方向相反）。\ce{BF3}：
非极性（平面内三个相同的\omterm{def:g11:polarity-and-cohesion:polar-bond}{极性键}成 \ang{120} 夹角，相互抵消）。
\end{solution}

\begin{solution}{exo:g11:polarity-and-cohesion:7}
\ce{O-H} 基团中的氢可以给出（不包括
与碳相连的 \ce{CH3} 中的氢）；氧\omterm{def:g7:atoms-and-molecules:atom}{原子}带两对\omterm{def:g10:lewis-and-shape:lone-pair}{孤电子对}，
可以接受。是的：甲醇的 \ce{O-H} 能与一个水\omterm{def:g7:atoms-and-molecules:molecule}{分子}的氧
成键，水的 \ce{O-H} 也能与甲醇的氧成键，例如
\ce{CH3-O-H}\,$\cdots$\,\ce{OH2}。
\end{solution}

\begin{solution}{exo:g11:polarity-and-cohesion:8}
三种\omterm{def:g7:atoms-and-molecules:molecule}{分子}都是非极性且形状相同，但越来越大
（分别含 10、18 和 36 个\omterm{def:g9:inside-the-atom:nucleus}{电子}）：范德华相互作用增强，
沸点也随之升高。
\end{solution}

\begin{solution}{exo:g11:polarity-and-cohesion:9}
第 14 族：甲烷中没有与 N、O 或 F 相连的氢，不形成
\omterm{def:g11:polarity-and-cohesion:hydrogen-bond}{氢键}，因此它遵循所在族的变化趋势。
\end{solution}

\begin{solution}{exo:g11:polarity-and-cohesion:10}
范德华相互作用：\omterm{def:g7:atoms-and-molecules:molecule}{分子}由 \ce{HCl} 增大到
\ce{HI}，这些相互作用随尺寸增大而增强的幅度
超过了\omterm{def:g11:polarity-and-cohesion:polar-bond}{部分电荷}之间吸引力的减小。
\end{solution}

\begin{solution}{exo:g11:polarity-and-cohesion:11}
碘\omterm{def:g7:atoms-and-molecules:molecule}{分子}远大于氯\omterm{def:g7:atoms-and-molecules:molecule}{分子}（\omterm{def:g9:inside-the-atom:nucleus}{电子}数 106
比 34）：其范德华相互作用要强得多，
强到足以使这些\omterm{def:g7:atoms-and-molecules:molecule}{分子}在室温下形成固体。
\end{solution}

\begin{solution}{exo:g11:polarity-and-cohesion:12}
丙烷是非极性的，只有范德华相互作用。甲氧基甲烷
为\omterm{def:g11:polarity-and-cohesion:polar-molecule}{极性分子}（两个极性 \ce{C-O} 键位于弯曲的 \ce{C-O-C} 中），但
其氧上没有氢，\omterm{def:g7:atoms-and-molecules:molecule}{分子}之间不形成\omterm{def:g11:polarity-and-cohesion:hydrogen-bond}{氢键}。乙醇
是\omterm{def:g11:polarity-and-cohesion:polar-molecule}{极性分子}，其 \ce{O-H} 基团可以形成\omterm{def:g11:polarity-and-cohesion:hydrogen-bond}{氢键}。吸引作用越强，
沸点越高：丙烷 < 甲氧基甲烷
< 乙醇。
\end{solution}

\begin{solution}{exo:g11:polarity-and-cohesion:13}
冰的开放网状结构比相同\omterm{def:g7:atoms-and-molecules:molecule}{分子}处于液体中
占据更多空间，因此冰的密度小于液态水，会漂浮。湖泊
自上而下凝固：冰层起隔热作用，使下方的水
仍保持液态，鱼因此得以存活。
\end{solution}

\begin{solution}{exo:g11:polarity-and-cohesion:14}
固体甲烷（小\omterm{def:g7:atoms-and-molecules:molecule}{分子}间的范德华相互作用）<
冰（\omterm{def:g11:polarity-and-cohesion:hydrogen-bond}{氢键}）<氯化钾（\omterm{def:g9:ions:ion}{离子}间的吸引）。
\end{solution}

\begin{solution}{exo:g11:polarity-and-cohesion:15}
一个水\omterm{def:g7:atoms-and-molecules:molecule}{分子}有两个氢\omterm{def:g7:atoms-and-molecules:atom}{原子}可以给出，两对\omterm{def:g10:lewis-and-shape:lone-pair}{孤电子对}可以接受，
即每个\omterm{def:g7:atoms-and-molecules:molecule}{分子}可参与四个\omterm{def:g11:polarity-and-cohesion:hydrogen-bond}{氢键}，其中两个给出、
两个接受，从而形成完整的网络。一个氟化氢\omterm{def:g7:atoms-and-molecules:molecule}{分子}
有三对\omterm{def:g10:lewis-and-shape:lone-pair}{孤电子对}，但只有一个氢\omterm{def:g7:atoms-and-molecules:atom}{原子}：它
只能给出一个\omterm{def:g11:polarity-and-cohesion:hydrogen-bond}{氢键}，因此平均每个\omterm{def:g7:atoms-and-molecules:molecule}{分子}只能形成一个
\omterm{def:g11:polarity-and-cohesion:hydrogen-bond}{氢键}（呈锯齿状链）。水\omterm{def:g7:atoms-and-molecules:molecule}{分子}平均形成的\omterm{def:g11:polarity-and-cohesion:hydrogen-bond}{氢键}数
约为其两倍，因此沸点更高。
\end{solution}

\begin{solution}{pb:g11:polarity-and-cohesion:1}
\textbf{1.} \ce{H-O-H}、\ce{H-S-H}、\ce{H-Se-H}，各中心\omterm{def:g7:atoms-and-molecules:atom}{原子}
均有两对\omterm{def:g10:lewis-and-shape:lone-pair}{孤电子对}：O、S 和 Se 的\omterm{def:g10:electron-shells:valence}{价电子}数均为 6，
它们同属一族。

\textbf{2.} $M(\ce{H2O}) = \qty{18.0}{g/mol}$，$M(\ce{H2S}) =
\qty{34.1}{g/mol}$，$M(\ce{H2Se}) = \qty{81.0}{g/mol}$。

\textbf{3.} $3.44 - 2.20 = 1.24$；$2.58 - 2.20 = 0.38$；
$2.55 - 2.20 = 0.35$。只有 \ce{O-H} 键明显是\omterm{def:g11:polarity-and-cohesion:polar-bond}{极性键}。

\textbf{4.} 折线形（四个价层\omterm{def:g9:inside-the-atom:nucleus}{电子}对，其中两对为\omterm{def:g10:lewis-and-shape:lone-pair}{孤电子对}）。水是\omterm{def:g11:polarity-and-cohesion:polar-molecule}{极性分子}；
硫化氢和硒化氢的极性则很微弱。

\textbf{5.} \omterm{def:g9:inside-the-atom:nucleus}{电子}数分别为 10、18 和 36：硒化氢
的范德华相互作用最强。

\textbf{6.} 只有水：其氢\omterm{def:g7:atoms-and-molecules:atom}{原子}与氧成键，而氧
是三种\omterm{def:g11:polarity-and-cohesion:electronegativity}{电负性}很强的\omterm{def:g7:atoms-and-molecules:atom}{原子}之一；硫和硒的\omterm{def:g11:polarity-and-cohesion:electronegativity}{电负性}
不够强，不足以使其氢\omterm{def:g7:atoms-and-molecules:atom}{原子}带上明显的
$\delta^+$。

\textbf{7.} 四个：两个通过自身的氢\omterm{def:g7:atoms-and-molecules:atom}{原子}提供，两个
由其两对\omterm{def:g10:lewis-and-shape:lone-pair}{孤电子对}接受。

\textbf{8.} 范德华相互作用（包括较弱的
微小\omterm{def:g11:polarity-and-cohesion:polar-bond}{部分电荷}之间的吸引）。

\textbf{9.} 分别为 $212.9 - 273.15 = \qty{-60.3}{\celsius}$ 和
\qty{-41.3}{\celsius}。

\textbf{10.} $-41.3 - (-60.3) = \qty{19.0}{\celsius}$。

\textbf{11.} 该\omterm{def:g7:atoms-and-molecules:molecule}{分子}较大，因此其范德华相互作用
更强。

\textbf{12.} $-60.3 - 19.0 = \qty{-79.3}{\celsius}$。

\textbf{13.} 升高量为 $-88.1 - (-112) = \qty{23.9}{\celsius}$；预测甲烷的
沸点为 $-112 - 23.9 = \qty{-136}{\celsius}$，与实际的 \qty{-162}{\celsius} 相比，
外推值约偏高 \qty{26}{\celsius}。这种外推
并不精确；真实的变化趋势并非直线。

\textbf{14.} 升高量为 $\qty{18.7}{\celsius}$；预测氟化氢的
沸点为 $-85.1 - 18.7 = \qty{-103.8}{\celsius}$，实际为 \qty{19.6}{\celsius}，
两者偏差约 \qty{123}{\celsius}。

\textbf{15.} $100 - (-79.3) = \qty{179}{\celsius}$，约
\qty{180}{\celsius}。

\textbf{16.} 升高量为 $\qty{25.3}{\celsius}$；预测沸点为 $-87.8 - 25.3 =
\qty{-113}{\celsius}$；偏差为 $-33.4 - (-113) = \qty{80}{\celsius}$。

\textbf{17.} 水（\qty{180}{\celsius}）> 氟化氢
（\qty{123}{\celsius}）> 氨（\qty{80}{\celsius}）。水既能提供两个
氢\omterm{def:g7:atoms-and-molecules:atom}{原子}，又以两对\omterm{def:g10:lewis-and-shape:lone-pair}{孤电子对}接受\omterm{def:g11:polarity-and-cohesion:hydrogen-bond}{氢键}，每个\omterm{def:g7:atoms-and-molecules:molecule}{分子}
共形成四个\omterm{def:g11:polarity-and-cohesion:hydrogen-bond}{氢键}，全部得到利用。氟化氢只有一个氢\omterm{def:g7:atoms-and-molecules:atom}{原子}可提供，而
氨只有一对\omterm{def:g10:lewis-and-shape:lone-pair}{孤电子对}可用于接受，平均每个\omterm{def:g7:atoms-and-molecules:molecule}{分子}
形成的\omterm{def:g11:polarity-and-cohesion:hydrogen-bond}{氢键}较少，且氮的\omterm{def:g11:polarity-and-cohesion:electronegativity}{电负性}较弱，所形成的\omterm{def:g11:polarity-and-cohesion:hydrogen-bond}{氢键}也较弱。

\textbf{18.} 它将是气体：其沸点接近 \qty{-80}{\celsius}。地球
上将没有液态水、海洋和降雨，也没有
依赖它们而存在的任何生命。

\textbf{19.} 它是对仅通过两点画出的一条直线进行外推，
而第 14 族的检验表明，此类外推的误差可达约
\qty{25}{\celsius}。

\textbf{20.} 若无\omterm{def:g11:polarity-and-cohesion:hydrogen-bond}{氢键}，水的沸点将约为
\qty{-80}{\celsius}，比其实际沸点低约
\qty{180}{\celsius}。
\end{solution}
